\documentclass[a4paper]{exam}
\usepackage[svgnames]{xcolor}
\usepackage{amsmath,amssymb,amsthm,titlesec,minted}

\newcommand
\lralign[2]{%
  \noindent\makebox[\textwidth]{\(\displaystyle#1\)\hfill\ifx&#2&\else(#2)\fi}\vspace{2ex}
}

\renewcommand{\thesubsection}{\alph{subsection}.}
\renewcommand{\thesubsubsection}{\roman{subsubsection}.}

\definecolor{faintgreen}{rgb}{0.5,1,0.5}
\definecolor{faintred}{rgb}{1,0.5,0.5}
\definecolor{faintblue}{rgb}{0.5,0.5,1}
\definecolor{faintyellow}{rgb}{1,1,0.5}
\definecolor{codebg}{RGB}{40,42,54}

\titleformat{\section}{\normalfont\Large\bfseries}{}{0pt}{}
\titleformat{\subsection}{\normalfont\large\bfseries}{}{0pt}{}
\titleformat{\subsubsection}{\normalfont\normalsize\bfseries\itshape}{}{0pt}{}
\titlespacing*{\section}{1em}{1.6\baselineskip}{0.6\baselineskip}
\titlespacing*{\subsection}{2em}{1.2\baselineskip}{0.4\baselineskip}
\titlespacing*{\subsubsection}{4em}{1.0\baselineskip}{0.3\baselineskip}

\setminted{style=dracula,bgcolor=codebg,breaklines,autogobble,fontsize=\small,tabsize=4,baselinestretch=1}
\NewDocumentCommand{\codefromfile}{O{}mm}{\par\vspace{0.4em}\inputminted[#1]{#2}{#3}\par\vspace{0.4em}}
\NewDocumentEnvironment{Section}{m}
{%
  \section*{#1}%
  \begingroup
  \leftskip=1em
}
{%
  \par
  \endgroup
}

\NewDocumentEnvironment{Subsection}{m}
{%
  \subsection*{#1}%
  \begingroup
  \leftskip=2em
}
{%
  \par
  \endgroup
}

\setlength{\parindent}{0in}
\setlength{\oddsidemargin}{0in}
\setlength{\textwidth}{6.5in}
\setlength{\textheight}{8.8in}
\setlength{\topmargin}{0in}
\setlength{\headheight}{18pt}

\printanswers
\title{Week 3}
\author{Radhesh Goel}
\date{\today}

\begin{document}
\maketitle

\section*{Question 1.}
\begin{solution}
  Fit a quadratic interpolating polynomial through the three measurements:
  \[
    (x_0,v_0)=(1.5,6.23),\quad (x_1,v_1)=(1.65,7.09),\quad (x_2,v_2)=(1.95,8.97)
  \]
  where $x$ is in metres and $v$ is in millimetres.

  The first divided differences are
  \[
    v[x_0,x_1] = \frac{7.09-6.23}{1.65-1.5} = 5.733333
  \]
  and
  \[
    v[x_1,x_2] = \frac{8.97-7.09}{1.95-1.65} = 6.266667
  \]
  Therefore, the Newton interpolating polynomial is
  \[
    v(x) = 6.23 + 5.733333(x - 1.5) + 1.185185(x - 1.5)(x - 1.65)
  \]
  Expanding:
  \[
    v(x) = 1.185185x^2+2x+0.563333
  \]
  with $v$ measured in millimetres.

  Differentiate twice:
  \[
    \begin{aligned}
      \frac{dv}{dx}     & = 2.370370x + 2 \\
      \frac{d^2v}{dx^2} & = 2.370370 \\
    \end{aligned}
  \]
  Because the interpolant is quadratic, the second derivative is constant.
  Therefore, at $x = 1.8\,\text{m}$,

  \[
    \boxed{
            \frac{d^2v}{dx^2} = 2.37037\ \text{mm}/\text{m}^2 = 2.37037 \times 10^{-3} \text{m}^{-1}
    }
  \]
\end{solution}

\section*{Question 2.}
\begin{solution}
  The beam equation is
  \[
    \frac{d^2v}{dx^2} = -\frac{M}{EI}
  \]
  Rearranging for $E$:
  \[
    E = -\frac{M}{I(d^2v/dx^2)}
  \]
  Converting the quantities to SI units:
  \[
    M = -21\ \text{kN} \cdot \text{m} = -21\,000\ \text{N} \cdot \text{m}
  \]
  and
  \[
    I = 18.5 \times 10^6\ \text{mm}^4.
  \]
  Since
  \[
    1\ \text{mm}^4=10^{-12}\ \text{m}^4
  \]
  then
  \[
    I = 18.5 \times 10^6 \times 10^{-12} = 1.85 \times 10^{-5}\ \text{m}^4
  \]

  Substitute into the equation:
  \[
    E = -\frac{-21\,000}{(1.85\times10^{-5})(2.37037\times10^{-3})}
  \]
  Therefore,
  \[
    E = 4.78885 \times 10^{11}\ \text{Pa}
  \]
  Thus, the estimated Young's modulus is
  \[
    \boxed{
	    E\approx4.79\times10^{11}\ \text{Pa} = 479\ \text{GPa}
	}
  \]
\end{solution}
\end{document}
